\subsection{叶状结构}

本节中，X 表示一个 C \(^{r}\) 类流形，其中 \(r \in \mathrm{N}_{\mathrm{K}}\)。除非另有说明，所考虑的所有流形和态射均假定属于 C \(^{r}\) 类。

9.2.1. 设 S 是一个流形，且 \(p: X \to S\) 是一个淹没。对每个 \(s \in S\)，赋予 \(X_s = p^{-1}(s)\) 由 X 的结构诱导的流形结构（5.10.5）；集合 X 是各 \(X_s\) 的不交并。记通过对 \(X_p\) 的 \(X_s\) 进行拼接而在 X 上得到的流形结构为 \(s \in S\)（5.2.4）；它是 X 上使各 \(X_s\) 成为开子流形的唯一流形结构。拓扑空间 \(X_p\) 是各空间 \(X_s\) 的和。

设 V 是一个流形，且 \(f\) 是从 V 到 X 的映射。要使 \(f\) 成为从 V 到 \(\mathbf{X}_{p}\) 的态射，充要条件是 \(f\) 是从 V 到 X 的态射，且 \(p \circ f\) 局部为常值。

9.2.2. X 的一个叶状结构是一个与 X 具有相同底层集合的流形 Y，并满足下列条件：

\begin{enumerate}
\def\labelenumi{(\Alph{enumi})}
\setcounter{enumi}{5}
\tightlist
\item
  对每个 \(x \in X\)，存在包含 x 的 X 的开子流形 U、一个流形 S 以及一个淹没 \(p: U \to S\)，使得流形 \(U_{p}\) 是 Y 的开子流形。
\end{enumerate}

也称对 (X, Y) 为叶状流形。若 (X, Y) 和 (X', Y') 是叶状流形，则从 (X, Y) 到 (X', Y') 的态射是任意映射 \(f: X \to X'\)，它既是从流形 X 到流形 X' 的态射，也是从流形 Y 到流形 Y' 的态射。

设 (X, Y) 是叶状流形。恒等映射 Y → X 是双射浸入。X 的子集 U 若是 Y 的开集，则称为一片叶；此时，赋予 U 由 Y 的拓扑结构和流形结构所诱导的拓扑结构和流形结构；例如，U 若是 Y 的开且连通子集，则称 U 为连通叶。作为 X 的子流形的叶构成 Y 的拓扑基。当 K = R 或 C 时，Y 的连通分支都是叶，称为 (X, Y) 的极大连通叶。

9.2.3. 若 \(s \in \mathbf{N}_{\mathbb{K}}\)、\(s \leqslant r\)，则将 \(\mathbf{C}^s\) 的 \(\mathbf{X}\) 类叶状结构称为以 \(\mathbf{C}^s\) 为底层的 \(\mathbf{X}\) 类流形的叶状结构。

9.2.4. 例子

\begin{enumerate}
\def\labelenumi{\alph{enumi})}
\item
  流形 X 是 X 的一个叶状结构，称为 X 的粗叶状结构。
\item
  若 \(p: \mathbf{X} \to \mathbf{S}\) 是淹没，则 \(\mathbf{X}_p\) 是 \(\mathbf{X}\) 的一个叶状结构，称为由 \(\mathbf{X}\) 定义的 \(p\) 的叶状结构。恒等映射定义的 \(\mathbf{X}\) 的叶状结构，是赋予了 0 维纯流形结构（5.2.1）的集合 \(\mathbf{X}\)；称为 \(\mathbf{X}\) 的离散叶状结构。
\item
  设 E 是 Banach 空间，F 是 E 的具有拓扑补空间的闭向量子空间，p 是典范投影 \(E \rightarrow E/F\)。E 的叶状结构 \(E_{p}\) 称为 F 定义的 E 的叶状结构。
\item
  设 \(\Gamma\) 是一个离散群，在拓扑空间 X 上适当地自由作用（TG, III, § 4, no. 4）。设 Y 是 X 的一个叶状结构，并假设对每个 \(s \in \Gamma\)，映射 \(x \mapsto sx\) 都是 (X, Y) 的自同构。在 X（分别在 Y）上以 \(\Gamma\) 的轨道为类的等价关系是正则的（5.9.5）。若将相应的商流形记为 X/\(\Gamma\)（分别为 Y/\(\Gamma\)），则 (X/\(\Gamma\),Y/\(\Gamma\)) 是一个叶状流形，称为 (X, Y) 按 \(\Gamma\) 的商。
\end{enumerate}

9.2.5. 设 Y 是 X 的一个叶状结构，U 是 X 的开子集。则 U 是 Y 中的开集，且赋予由 Y 诱导的流形结构后的 U 是 X 的一个叶状结构，称为由 Y 诱导的叶状结构。

更一般地，设 \(f: X' \to X\) 是一个态射，使得 \(f\) 与恒等映射 \(Y \to X\) 构成横截对（5.11.1）。纤维积 \(Y' = X' \times_X Y\) 典范地等同于 \(X'\) 的一个叶状结构，称为 \(Y\) 沿 \(f\) 的拉回叶状结构。

9.2.6. 设 \(p: \mathbf{X} \to \mathbf{S}\) 和 \(p': \mathbf{X}' \to \mathbf{S}'\) 是两个淹没，且 \(f\) 是从 \(\mathbf{X}\) 到 \(\mathbf{X}'\) 的态射。要使 \(f\) 成为从 \(\mathbf{X}_p\) 到 \(\mathbf{X}_{p'}\) 的态射，充要条件是：对于每个 \(x \in \mathbf{X}\)，存在一个开邻域 \(\mathbf{U}\)，它是 \(x\) 在 \(\mathbf{X}\) 中的开邻域，以及一个态射 \(g\)，定义在 \(p(\mathbf{U})\) 上且取值于 \(\mathbf{S}'\)，使得 \(g \circ p = p' \circ f\) 在 \(\mathbf{U}\) 上成立。

9.2.7. 设 Y 是 X 的一个叶状结构。(X, Y) 的一个叶状坐标图是四元组 (U, \(\varphi\), E, F)，其中 (U, \(\varphi\), E) 是 X 的坐标图，F 是 Banach 空间 E 的闭向量子空间，具有拓扑补空间，并且 \(\varphi\) 是从由 Y 诱导的 U 的叶状结构到由从 E 到 E/F 的典范映射定义的 \(\varphi(\mathbf{U})\) 的叶状结构的同构。对于每个点 \(x \in \mathbf{X}\)，存在 (X, Y) 的叶状坐标图 (U, \(\varphi\), E, F)，使得 \(x \in \mathbf{U}\)。令 \(n = \dim_x \mathbf{X}\)，\(m = \dim_x \mathbf{Y}\)。若 \(n\) 有限，则存在 \(x\) 的一个开邻域 U 以及 X 在 U 上的坐标系 \(\zeta^1, \ldots, \zeta^n\)，使得由 Y 诱导的 U 的叶状结构与由映射（\(\zeta^{m+1}, \ldots, \zeta^n\)）从 U 到 \(\mathbf{K}^{n-m}\) 所定义的叶状结构一致。

9.2.8. 设 Y 是 X 的一个叶状结构。当 x 遍历 X 时，空间 \(T_{x}(Y)\) 是 T(X) 的一个 \(C^{r-1}\) 类向量子丛的纤维；该子丛称为与叶状结构 Y 相切的 T(X) 的子丛，记为 T(X, Y)。若叶状结构 Y 由淹没 \(p: X \to S\) 定义，则 T(X, Y) = Ker T(p) 是 X 关于 S 的相对切丛 T(X/S)（8.1.3）。

设 (X, Y) 和 (X', Y') 是两个叶状流形，且 f 是从 X 到 X' 的态射。f 是叶状流形的态射的充要条件是 T(f) 将 T(X, Y) 映入 T(X', Y')。特别地，设 \(f: Z \to X\) 是流形的态射；f 是从 Z 到 Y 的态射的充要条件是 T(f) 将 T(Z) 映入 T(X, Y)。X 的子流形 Z 是 (X, Y) 的一片叶的充要条件是有

\[
\mathrm{T} _ {z} (\mathrm{Z}) = \mathrm{T} _ {z} (\mathrm{X}, \mathrm{Y}) \quad \text {for all} \quad z \in \mathrm{Z}.
\]

设 Y 是 X 的一个叶状结构，U 是 Y 的一片叶，并具有可数的开集基。设 Z 是一个流形，f 是从 Z 到 U 的映射。f 是从 Z 到 U 的流形态射的充要条件是 f 与 U 到 X 的典范嵌入的复合是从 Z 到 X 的流形态射。若 X 局部紧且在无穷远处可数，则 (X, Y) 的每片连通叶都具有可数的开集基（参见~TG, I, 第 3 版，§11，第 7 号，定理 1 的推论 2）。

9.2.9. 假设 K = R 或 C。设 Y 是 X 的一个叶状结构。以下条件等价：

\begin{enumerate}
\def\labelenumi{\alph{enumi})}
\item
  存在一个流形 S 和一个淹没 \(p: X \to S\)，使得 \(Y = X_p\)。
\item
  对每个 \(x \in X\)，存在 X 的一个子流形 \(S_{x}\)，具有下列两个性质：
\begin{enumerate}
\item[\(b_1)\)] 切空间 \(\mathbf{T}_{x}(\mathbf{S}_{x})\) 是 \(\mathbf{S}_{x}\) 在 \(x\) 处的切空间，是 \(\mathbf{T}_{x}(\mathbf{X},\mathbf{Y})\) 中 \(\mathbf{T}_{x}(\mathbf{X})\) 的拓扑补空间。
\item[\(b_2)\)] (X, Y) 的每片连通叶与 \(S_{x}\) 至多相交于一点。
\end{enumerate}
\end{enumerate}

假设这些条件满足，令 R 为 X 上的关系“x 和 y 属于同一片连通叶”；则 R 是 X 上的正则等价关系（5.9.5），并且若 p 表示典范投影 \(X \to X/R\)，则有 \(Y = X_p\)。

例。--- 取 K = R 和 X = R²。令离散群 Γ = Z² 通过平移作用于 X。令 m ∈ R，令 (X, Yₘ) 为由淹没 p: (x₁, x₂) ↦ x₂ - mx₁ 从 X 到 R 所定义的叶状结构。置 X' = X/Γ 和 Yₘ' = Yₘ/Γ；则 Yₘ' 是环面 X' 的一个叶状结构（参见~9.2.4，例 d)）。若 m 是有理数，则存在从 X' 到 R/Z 的淹没 p'，使得 Yₘ' = Xₚ'。若 m 是无理数，则 (X', Yₘ') 的每片极大连通叶在 X' 中稠密，并且不存在从 X' 到流形 S 的淹没 p'，使得 Yₘ' = Xₚ'。

9.2.10. 设 Y 是 X 的一个叶状结构，且 \(\pi: X \to S\) 是流形的态射，使得复合 \(Y \to X \xrightarrow{\pi} S\) 是平展的。若 \(S'\) 是 S 的开子集，则称截面 \(\sigma: S' \to X\) 是 \(\pi\) 在 \(S'\) 上的截面；相对于叶状结构 Y，若它是从 \(S'\) 到 Y 的态射，则称其为水平截面；这等价于说 \(\sigma(S')\) 是 (X, Y) 的一片叶，或者说 \(T(\sigma)\) 将 \(T(S')\) 映到 \(T(X, Y)\)。对于每个 \(s_0 \in S\) 以及每个 \(x_0 \in \pi^{-1}(s_0)\)，存在一个定义在 \(s_0\) 的邻域内且在 \(x_0\) 处取值 \(s_0\) 的水平截面；两个这样的截面在 \(s_0\) 的某个邻域内一致。

更一般地，设 \(s_{0}\in S\)，T 是一个流形，设 \(f\) 是从 T 到 X 的子流形 \(\pi^{-1}(s_{0})\) 的态射，且 \(t_{0}\in T\)。则存在一个开邻域 \(\mathbf{T}^{\prime}\)（分别为一个开邻域 \(S'\)），其中前者包含 \(t_{0}\)，后者包含 \(s_{0}\)，以及一个态射

\[
\mathbf {F}: \mathbf {S} ^ {\prime} \times \mathbf {T} ^ {\prime} \rightarrow \mathbf {X}
\]

满足下列条件：对于每个 \(t \in \mathbf{T}'\)，映射 \(s \mapsto \mathrm{F}(s, t)\) 是 \(\pi\) 在 \(\mathbf{S}'\) 上的水平截面，并取值为 \(f(t)\)，取值点为 \(s_0\)。两个这样的映射 F 在 \((s_0, t_0)\) 的某个邻域内一致。

